%! Least Squares Approximations — Unit 4, Lesson 4.3 — Student Packet Cover Sheet
\documentclass[10pt]{article}
\usepackage{linalg-article}
\usepackage{linalg-boxes}
\usepackage{ltablex}
\keepXColumns

\begin{document}

% Full-bleed burgundy banner
\begin{tikzpicture}[remember picture, overlay]
  \fill[burgundy] (current page.north west)
    rectangle ([yshift=-0.9in]current page.north east);
\end{tikzpicture}
\vspace{-0.8in}

\begin{center}
  {\color{white}\textbf{\LARGE Linear Algebra}}\\[4pt]
  {\color{blushmid}\large\textbf{Unit 4 \quad Orthogonality}}\\[6pt]
  {\color{white}\textbf{\large Lesson 4.3 \quad Least Squares Approximations}}
\end{center}

\vspace{0.22in}
\namedateperiod
\vspace{0.18in}

\begin{learningtargetbox}
\small
\begin{enumerate}[leftmargin=1.5em, itemsep=5pt]
  \item \ldots{} explain why $A\mathbf{x}=\mathbf{b}$ can have \textbf{no} solution, and what a \textbf{least-squares} best answer means.
  \item \ldots{} fit a \textbf{best-fit line} $y=C+Dt$ to data by solving the \textbf{normal equations} $A^{\mathsf{T}}A\hat{\mathbf{x}}=A^{\mathsf{T}}\mathbf{b}$.
  \item \ldots{} find the fitted values and \textbf{residuals}, and explain \emph{why} least squares makes the total squared error smallest.
\end{enumerate}
\end{learningtargetbox}

\vspace{0.14in}

\begin{tocbox}
\small
\renewcommand{\arraystretch}{1.5}
\begin{tabularx}{\linewidth}{@{} c l X r @{}}
\rowcolor{blush}
\textbf{\#} & \textbf{Component} & \textbf{Description} & \textbf{Score} \\
\midrule
1 & Warm-Up        & Spiral review: reading a line \& residual, building $A^{\mathsf{T}}A$, solving a $2\times2$ system & \blank{1.2cm} \\
2 & Guided Notes   & No-solution systems, best-fit line, normal equations, fitted values \& residuals & \blank{1.2cm} \\
3 & Group Activity & Best Fit by Least Squares --- a flat line, a line through 3 points, and a prediction & \blank{1.2cm} \\
4 & Exit Ticket    & Quick check: fit a line, read the residual, and why not solve $A\mathbf{x}=\mathbf{b}$ & \blank{1.2cm} \\
5 & Homework       & Best constant, line fits, prediction, and where the residual lives & \blank{1.2cm} \\
\midrule
  & \multicolumn{2}{r}{\textbf{Total}} & \blank{1.2cm} \\
\end{tabularx}
\end{tocbox}

\vspace{0.10in}

\begin{remindbox}
\small One idea drives the lesson: when data will not sit on a line, $A\mathbf{x}=\mathbf{b}$ has
\emph{no} solution, so we take the \textbf{closest} reachable point $\mathbf{p}=A\hat{\mathbf{x}}$ --- the
\textbf{projection} of the data (Lesson~4.2). Solving the \textbf{normal equations}
$A^{\mathsf{T}}A\hat{\mathbf{x}}=A^{\mathsf{T}}\mathbf{b}$ gives the best-fit line, and the leftover
\textbf{residual} $\mathbf{e}=\mathbf{b}-\mathbf{p}$ --- the part least squares makes small --- lands in
the orthogonal complement $N(A^{\mathsf{T}})$ (Lesson~4.1).
\end{remindbox}

\end{document}
